\centerline{\bf \Huge Финальная МатДрака}

\textit{ВСЕ участники команды, занявшей I место, получают по 8 баллов в рейтинговую таблицу и право на 8 попыток ввода кода.}

\textit{ВСЕ участники команды, занявшей II место, получают по 4 балла в рейтинговую таблицу и право на 4 попытки ввода кода.}

\textit{ВСЕ участники команды, занявшей III место, ничего не получают.}

\textbf{Ученик, решивший задачу, получает 6 баллов в рейтинговую таблицу и право на 6 попыток ввода кода; если решали двое, то по 3 балла и право на 3 попытки, если решали трое, то по 2 балла и право на 2 попытки}\\

\problem{\textbf{1}} Является ли число 102030405060708090807060504030201 квадратом какого-нибудь целого числа?\\


\problem{\textbf{2}} На какое целое число надо умножить 999 999 999, чтобы получить число, состоящее из одних единиц?\\

\problem{\textbf{3}} Найдите наибольшее натуральное число, каждая некрайняя цифра которого меньше среднего арифметического соседних с ней цифр.\\

\problem{\textbf{4}} В левом нижнем углу клетчатой доски $n\times n$ стоит конь. Известно, что наименьшее число ходов, за которое конь может дойти до правого верхнего угла, равно наименьшему числу ходов, за которое он может дойти до правого нижнего угла. Найдите все такие $n$.\\

\problem{\textbf{5}} В футбольном турнире участвовало 8 команд, причём каждая сыграла с каждой ровно по одному разу. Известно, что каждые две команды, сыгравшие между собой вничью, набрали в итоге разное число очков. Найдите наибольшее возможное общее число ничьих в этом турнире. (За выигрыш матча команде начисляется 3 очка, за ничью – 1, за поражение – 0.)\\




\problem{\textbf{6}} Каждая клетка таблицы размером $7\times8$ (7 строк и 8 столбцов) покрашена в один из трёх цветов: красный, жёлтый или зелёный. При этом в каждой строке красных клеток не меньше, чем жёлтых, и не меньше, чем зелёных, а в каждом столбце жёлтых клеток не меньше, чем красных, и не меньше, чем зелёных. Сколько зелёных клеток может быть в такой таблице?\\






\problem{\textbf{7}} Шахматная фигура может сдвигаться на 8 или 9 клеток по горизонтали или вертикали. Запрещается ходить на одну и ту же клетку дважды.
Какое наибольшее количество клеток может обойти эта фигура на доске $15\times15$? (Начать можно с любой клетки.)\\



\problem{\textbf{8}} В таблице $2005\times2006$ расставлены числа 0, 1, 2 так, что сумма чисел в каждом столбце и в каждой строке делится на 3.
Какое наибольшее возможное количество единиц может быть в этой таблице?\\


\problem{\textbf{9}}  На книжной полке стоят 30 томов энциклопедии в некотором порядке. За одну операцию разрешается менять местами любые два соседних тома. За какое наименьшее число операций можно гарантированно выстроить все тома в правильном порядке (с первого по тридцатый слева направо) независимо от начального положения?\\

\problem{\textbf{10}} Найдите $\dfrac{X}{100}$, если:

$\dfrac{1}{1 \cdot 2 \cdot 3}+\dfrac{1}{2 \cdot 3 \cdot 4}+\dfrac{1}{3 \cdot 4 \cdot 5}+\dfrac{1}{4 \cdot 5 \cdot 6}+\ldots+\dfrac{1}{2025 \cdot 2026 \cdot 2027}=\dfrac{X}{4\cdot2026\cdot2027}$.\\

\problem{\textbf{11}} Можно ли расставить натуральные числа от 1 до 10 в ряд так, чтобы каждое число было делителем суммы всех предыдущих?\\

\problem{\textbf{12}} Найдите ВСЕ решения ребуса: МАД + РИД = РЕАЛ\\

\problem{\textbf{13}} В треугольнике $AIM$ угол $AIM=100^{\circ}$. Точка $T$ лежит на стороне $AM$ так, что $IT$ является биссектрисой треугольника $AIM$. Найдите угол $AMI$, если $IT=AT$.





